% Part: first-order-logic
% Chapter: beyond
% Section: modal-logics

\documentclass[../../../include/open-logic-section]{subfiles}

\begin{document}

\olfileid{fol}{byd}{mod}

\olsection{منطق‌های وجهی}

مثال زیر را از یک جملهٔ شرطی در نظر بگیرید:
\begin{quote}
  اگر جرمی در آن اتاق تنها باشد، پس مست و برهنه است و روی صندلی‌ها
  می‌رقصد.
\end{quote}
این نمونه‌ای از گزاره‌ای شرطی است که ممکن است به لحاظ مادی صادق اما
بااین‌همه گمراه‌کننده باشد، زیرا گویی القا می‌کند میان مقدم و نتیجه
پیوندی نیرومندتر از این برقرار است که صرفاً یا مقدم کاذب باشد یا تالی
صادق. یعنی عبارت‌بندی القا می‌کند که ادعا نه‌تنها در همین جهان صادق است
(جایی که ممکن است به‌سادگی از آن رو صادق باشد که جرمی در اتاق تنها
نیست)، بلکه افزون بر آن، اگر مقدم صادق emph{می‌بود}، نتیجه نیز صادق
\emph{می‌بود}. به بیان دیگر، می‌توان مراد از این گزاره را آن دانست که
ادعا نه‌فقط در این جهان، بلکه در هر جهان «ممکن» صادق است؛ یا اینکه در
برابر صرفاً صادق‌بودن در همین جهان خاص، emph{ضرورتاً} صادق است.

منطق وجهی برای معناکردن این‌گونه ضرورت طراحی شد. منطق گزاره‌ای وجهی با
افزودن یک عملگر جعبه به منطق گزاره‌ای عادی به دست می‌آید؛ یعنی اگر
$!A$ یک !!a{formula} باشد، $\Box !A$ نیز چنین است. از دیدگاه شهودی،
$\Box !A$ می‌گوید $!A$ emph{ضرورتاً} صادق، یا در هر جهان ممکنی صادق
است. معمولاً $\Diamond !A$ را اختصاری برای $\lnot \Box \lnot !A$
می‌گیرند و می‌توان آن را چنین خواند که $!A$ emph{ممکن است} صادق باشد.
البته وجهیت را می‌توان به منطق محمولات نیز افزود.

می‌توان از !!{structure}s کریپکی برای فراهم‌کردن معناشناسی منطق وجهی
استفاده کرد؛ در واقع کریپکی نخست این معناشناسی را با نظر به منطق وجهی
طراحی کرد. به‌جای محدودشدن به ترتیب‌های جزئی، در حالت کلی‌تر مجموعه‌ای
از «جهان‌های ممکن» چون~$P$ و رابطهٔ دوتایی «دسترسی»
$\Atom{R}{x,y}$ میان جهان‌ها داریم. از دیدگاه شهودی، $\Atom{R}{p,q}$
می‌گوید جهان~$q$ با~$p$ سازگار است؛ یعنی اگر «در» جهان~$p$ باشیم، باید
این امکان را جدی بگیریم که جهان می‌توانست مانند~$q$ باشد.

منطق وجهی را گاهی در برابر منطقی «مصداقی»، منطقی «شدتی» می‌نامند.
معناشناسی مقصود برای منطقی مصداقی، مانند منطق کلاسیک، تنها به یک جهان،
یعنی جهان «واقعی»، ارجاع می‌دهد؛ حال آنکه معناشناسی منطقی «شدتی» بر
هستی‌شناسی پرمایه‌تری تکیه دارد. افزون بر !!{structure}دادن به ضرورت،
می‌توان وجهیت را برای !!{structure}دادن به دیگر ساخت‌های زبانی نیز به
کار برد و $\Box$ و $\Diamond$ را متناسب با کاربرد از نو تفسیر کرد.
برای نمونه:
\begin{enumerate}
\item در منطق اثبات‌پذیری، $\Box !A$ چنین خوانده می‌شود: «$!A$
  اثبات‌پذیر است»، و $\Diamond !A$ چنین: «$!A$ سازگار است.»
\item در منطق معرفتی، می‌توان $\Box !A$ را چنین خواند: «می‌دانم
  $!A$» یا «باور دارم $!A$.»
\item در منطق زمانی، می‌توان $\Box !A$ را چنین خواند: «$!A$ همیشه
  صادق است»، و $\Diamond !A$ را چنین: «$!A$ گاهی صادق است.»
\end{enumerate}

می‌خواهیم منطق را با قواعد و اصول موضوعی که با وجهیت سروکار دارند
گسترش دهیم. برای نمونه، دستگاه~$\Log{S4}$ از اصول موضوع و قواعد عادی
منطق گزاره‌ای، همراه با اصول موضوع زیر، تشکیل می‌شود:
\begin{align*}
& \Box (!A \lif !B) \lif (\Box !A \lif \Box !B)\\
& \Box !A \lif !A \\
& \Box !A \lif \Box \Box !A
\intertext{و نیز قاعده‌ای که می‌گوید «از $!A$، $\Box !A$ را نتیجه
  بگیرید.» $\Log{S5}$ اصل موضوع زیر را می‌افزاید:}
& \Diamond !A \lif \Box \Diamond !A
\end{align*}
صورت‌های دیگری از این اصول موضوع ممکن است برای کاربردهای متفاوت مناسب
باشند؛ برای نمونه، معمولاً S5 را مشخصه‌ساز مفهوم ضرورت منطقی می‌گیرند.
نکتهٔ خوشایند آن است که غالباً می‌توان معناشناسی‌ای یافت که دستگاه
!!{derivation} نسبت به آن صحیح و تام باشد؛ کافی است رابطهٔ دسترسی در
!!{structure}s کریپکی را به شیوه‌هایی طبیعی محدود کنیم. برای نمونه،
$\Log{S4}$ با ردهٔ !!{structure}s کریپکی‌ای متناظر است که رابطهٔ دسترسی
در آن‌ها بازتابی و تراگذری است. $\Log{S5}$ با ردهٔ !!{structure}s
کریپکی‌ای متناظر است که رابطهٔ دسترسی در آن‌ها {\em همگانی} است؛ یعنی
هر جهان از هر جهان دیگر دسترس‌پذیر است. پس $\Box !A$ هرگاه و تنها
هرگاه برقرار است که $!A$ در هر جهانی برقرار باشد.

\end{document}
